% appendix_traits.tex: is the rate a characteristic of the person or a product of circumstances? (memo traits_vs_circumstances); moved from Section 10.2 with Table tab:person_place.
\section{A Characteristic of the Person or a Product of Circumstances?}\label{appendix:traits}

Write $r(t, \iota)$ for the discount rate at horizon $t$ of a household of type $\iota$, the type being whatever the score, income, place and cohort of Section \ref{sec:heterogeneity} index. The schedule is written as if it belonged to the household. Whether it does is an empirical question with a literature on each side, and the answer determines what the heterogeneity in Section \ref{sec:heterogeneity} means. Measured discount rates are moderately stable within a person over a year \citep{meier_temporal_2015}, and a single elicitation predicts wealth decades later \citep{epper_time_2020}, present bias predicts how card debt is paid down \citep{kuchler_sticking_2021}, and impatience measured in adolescence predicts lifetime earnings and credit outcomes \citep{cadena_human_2015}. Financial distress is highly persistent within the person, and a model reproduces that persistence only with persistent heterogeneity in discount factors \citep{athreya_persistence_2019}; the hand-to-mouth are, on this view, impatient by preference rather than constrained by circumstance \citep{aguiar_who_2024}. Against this, the modern evidence on default is a literature about circumstances: nearly every mortgage default is preceded by a shortfall in cash flow \citep{ganong_why_2023}, most defaulters cannot pay rather than choose not to \citep{gerardi_cant_2018}, and more generous unemployment insurance lowers default \citep{hsu_unemployment_2018}. The two literatures answer different questions. Circumstances govern the timing of a default; persistent characteristics govern who defaults and who defaults repeatedly: fewer than a tenth of borrowers account for half of all severe-delinquency years \citep{athreya_persistence_2019}, and the variance of the marginal propensity to consume splits roughly in half between a persistent household component and temporary cash on hand \citep{gelman_what_2021}. The one design that treats person and place symmetrically follows borrowers who move: \citet{keys_place_2023} find that adult moves close less than one-tenth of the geographic gap in collections and card non-payment, so that the person accounts for the remaining nine-tenths, while moves close a quarter to a third of the gap in bankruptcy filing, which depends on local information and lawyers. The person component they measure is portable rather than necessarily preference-based, and childhood exposure to place forms much of it. The weight of the evidence is therefore that whether a household defaults is mostly a characteristic that moves with it, that when it defaults is mostly a matter of circumstance, and that the two interact.

This paper's estimates bear on the question in four places. First, the payment margin is unchanged by borrower fixed effects: among borrowers with two or more auto loans, the same person defaults more on the contract with the larger payment (Table \ref{tab:main_margins}, lower panel), so the response to the contract is not an artifact of who chooses which contract. Second, distress is persistent but not absorbing: 20.5 percent of delinquent auto trades cure within a four-month archive gap (Appendix \ref{appendix:measuredrho}), and the rating-transition matrix supports a decomposition into chronic and transitory delinquency that the marginal-utility correction of Appendix \ref{appendix:mucorrection} uses. Third, the gradient of the payment margin in the credit score is a stable characteristic modulating the response to a circumstance, since the score is the market's estimate of the person's persistent component. Fourth, the payment margin is estimated separately for cohorts that faced different local labor-market shocks in their first eighteen months (Section \ref{sec:liquidity}); if the coefficient were a repackaged liquidity shock it would move with the shock.

Two designs sharpen the decomposition on the credit-bureau data. Among repeat borrowers, the covariance of a borrower's default residuals across two of her own loans, after lender-year and zip3-year effects are removed, estimates the variance of the persistent borrower component without a degrees-of-freedom correction, and its share of the residual variance is the person's share of default; its correlation with the score says how much of that share the market prices. And borrowers whose zip3 changes between loans identify the place component in the manner of \citet{chetty_impacts_2018} and \citet{keys_place_2023}: the change in a mover's own default across the two loans, regressed on the change in the stayers' mean default between destination and origin, has a slope of one if the geographic gradient is causal and zero if it is composition. Table \ref{tab:person_place} reports both. Their implication for the paper is direct: the geographic heterogeneity in Section \ref{sec:heterogeneity} is a property of places only to the extent of the movers' slope, $0.075$ ($0.013$) over the life of the loan and $0.119$ ($0.017$) within 24 months, and the type axis $\iota$ mixes persistent characteristics (the score) with circumstances (income at origination) in proportions those two exhibits measure.


\paragraph{Place and the vote.} If partisan differences in attitudes toward indebtedness were differences in how far payments are weighed, the far-payment weight would show a gradient in the presidential vote across states. Figure \ref{fig:geo_politics} (Appendix \ref{appendix:figures}) relates the payment elasticity, the default rate and the term-to-payment ratio to the vote by state. The Republican share is the two-party share from county returns summed to the state \citep{mcgovern_county_returns}, and the 2016 shift is the 2016 share minus the 2012 share, which runs from $-0.12$ in Utah to $0.09$ in West Virginia; the sample is the 43 states with at least 300 charge-offs within 24 months and a vote share. The payment elasticity is lower where the Republican share is higher: the loans-weighted slope is $-0.67$ (heteroskedasticity-robust standard error $0.15$), a weighted correlation of $-0.54$, and it is $-0.55$ ($0.19$) without the five largest states and $-0.49$ ($0.16$) within census region. The default rate is higher, $0.022$ ($0.009$) per unit of the share on a mean of $0.019$. Both gradients are composition: the share is correlated $-0.73$ with mean income across states, and with the state's mean score, income and contract length controlled the elasticity's slope is $-0.13$ ($0.24$) and the default rate's changes sign. The far payments' weight shows no gradient. The ratio of the term elasticity to the payment elasticity has a slope of $-0.43$ ($0.50$) on the share, a correlation of $-0.12$, and after the adjustments of Appendix \ref{appendix:ladder} the rank correlation of the rate the formula returns with the share is $-0.02$; the ratio's sign on the share is not stable across outcomes. Neither the exposure objects nor the ratio relate to the shift: the slopes are $0.08$ ($0.53$) for the elasticity, $-0.04$ ($0.03$) for the default rate and $0.78$ ($1.22$) for the ratio. Places that vote Republican hold more households near the default threshold and their default responds less to the payment; they weigh far payments as everyone else does.

\begin{table}[ht]
\centering
\caption{Person and place in auto-loan default: variance decomposition and movers}
\label{tab:person_place}
\resizebox{\textwidth}{!}{\begin{tabular}{lrrrr}
\toprule
\multicolumn{5}{l}{\emph{Panel A: the person's share of the variance of default (two-loan covariance, repeat borrowers)}} \\
Outcome & loans & borrowers & person share & corr.\ with score \\
\midrule
over the life & 64,357,272 & 13,521,907 & 0.245 & -0.250 \\
within 24 months & 59,133,648 & 13,417,298 & 0.269 & -0.158 \\
within 12 months & 62,496,464 & 13,509,178 & 0.281 & -0.103 \\
\midrule
\multicolumn{5}{l}{\emph{Panel B: movers -- slope of the change in own default on the change in the destination's same-year default rate}} \\
Outcome & movers & & slope & (s.e.) \\
\midrule
over the life, no controls & 1,841,048 & & 0.075 & (0.013) \\
within 24 months, no controls & 1,787,142 & & 0.119 & (0.017) \\
over the life, with controls & 1,841,048 & & 0.073 & (0.013) \\
within 24 months, with controls & 1,787,142 & & 0.119 & (0.017) \\
\bottomrule
\end{tabular}
}
% replication: Code/identification/I7_person_place.R, Code/exhibits/make_placeholder_tables.R
\end{table}
